\section*{Canonical eutectic path and liquor interpolation}

We consider a canonical eutectic path discretised into nodes and segments.

\subsection*{Nodes and solids from gradients/intercepts}

Each canonical node \(m\) has an associated eutectic composition


\[
\mathbf{X}^{(m)}_{\mathrm{eu}} = \bigl(X^{(m)}_{\mathrm{eu},1}, \dots, X^{(m)}_{\mathrm{eu},N}\bigr),
\]


obtained from the eutectic model.

Along the path, the solid compositions at the nodes are not arbitrary;
they are determined by the canonical gradients and intercepts of the
eutectic path. For each component \(k\),


\[
X^{(m)}_{\mathrm{s},k}
= a_k + b_k \, f_m,
\]


where \(a_k\) and \(b_k\) are the intercept and gradient for component \(k\),
and \(f_m\) is the canonical path parameter at node \(m\).
Thus the solid compositions


\[
\mathbf{X}^{(m)}_{\mathrm{s}} = \bigl(X^{(m)}_{\mathrm{s},1}, \dots, X^{(m)}_{\mathrm{s},N}\bigr)
\]


are fully determined by the gradient/intercept representation of the
eutectic path. These are the \(X_s\) values used in the slicing engine.

\subsection*{Segments and recovery intervals}

The path is partitioned into segments indexed by \(j\).
Each segment \(j\) covers a recovery interval


\[
R^{(j)}_{\mathrm{start}} \quad \text{to} \quad R^{(j)}_{\mathrm{end}},
\]


with an associated solid composition \(\mathbf{X}^{(j)}_{\mathrm{s}}\)
(which is constant within the segment in the canonical representation).

For each segment we also have the liquor compositions at the segment
boundaries:


\[
\mathbf{X}^{(j)}_{\mathrm{liq,start}}
\quad \text{and} \quad
\mathbf{X}^{(j)}_{\mathrm{liq,end}},
\]


which encode the cumulative effect of all previous eutectic removals
up to \(R^{(j)}_{\mathrm{start}}\) and \(R^{(j)}_{\mathrm{end}}\),
respectively. In other words, the history of “subtracting eutectic”
is embedded in these boundary liquors; we do not recompute the full
liquor path at slice time.

\subsection*{Liquor interpolation within a segment}

Given a target recovery \(R\) lying inside segment \(j\),


\[
R^{(j)}_{\mathrm{end}} \le R \le R^{(j)}_{\mathrm{start}},
\]


we define a local interpolation parameter


\[
t(R) =
\frac{R^{(j)}_{\mathrm{start}} - R}
     {R^{(j)}_{\mathrm{start}} - R^{(j)}_{\mathrm{end}}}
\quad \in [0,1].
\]



The liquor composition at recovery \(R\) is then obtained by linear
interpolation between the segment boundary liquors:


\[
\mathbf{X}_{\mathrm{liq}}(R)
= (1 - t(R)) \, \mathbf{X}^{(j)}_{\mathrm{liq,start}}
  + t(R) \, \mathbf{X}^{(j)}_{\mathrm{liq,end}}.
\]



Componentwise, for each \(k = 1, \dots, N\),


\[
X_{\mathrm{liq},k}(R)
= (1 - t(R)) \, X^{(j)}_{\mathrm{liq,start},k}
  + t(R) \, X^{(j)}_{\mathrm{liq,end},k}.
\]



Because both \(\mathbf{X}^{(j)}_{\mathrm{liq,start}}\) and
\(\mathbf{X}^{(j)}_{\mathrm{liq,end}}\) are valid compositions
(i.e. each component lies in \([0,1]\) and the vectors are normalised),
the interpolated liquor \(\mathbf{X}_{\mathrm{liq}}(R)\) is also
bounded and normalised.

\subsection*{Consistency with the lever rule}

At any slice recovery \(R\), the feed composition \(\mathbf{X}_{\mathrm{c}}\),
solid composition \(\mathbf{X}_{\mathrm{s}}\), and liquor composition
\(\mathbf{X}_{\mathrm{liq}}(R)\) satisfy the lever rule:


\[
\mathbf{X}_{\mathrm{c}}
= R \, \mathbf{X}_{\mathrm{s}}
  + (1 - R) \, \mathbf{X}_{\mathrm{liq}}(R),
\]


or componentwise,


\[
X_{\mathrm{c},k}
= R \, X_{\mathrm{s},k}
  + (1 - R) \, X_{\mathrm{liq},k}(R).
\]



Equivalently, for each component \(k\),


\[
R
= \frac{X_{\mathrm{c},k} - X_{\mathrm{liq},k}(R)}
       {X_{\mathrm{s},k} - X_{\mathrm{liq},k}(R)},
\]


and all components yield the same recovery \(R\) (within numerical
tolerance). This is exactly the mass-balance identity you are checking
in the Excel output.

\subsection*{Summary}

- The solid compositions \(\mathbf{X}_{\mathrm{s}}\) used in the slices
  come directly from the canonical gradients and intercepts of the
  eutectic path; they are not invented at slice time.
- The liquor compositions along the path are represented by their
  values at segment boundaries, \(\mathbf{X}^{(j)}_{\mathrm{liq,start}}\)
  and \(\mathbf{X}^{(j)}_{\mathrm{liq,end}}\), which encode the cumulative
  effect of eutectic removal.
- For any slice recovery \(R\) inside a segment, the liquor composition
  is obtained by linear interpolation between these boundary liquors,
  ensuring bounded, normalised, and mass-balance-consistent
  \(\mathbf{X}_{\mathrm{liq}}(R)\).
